\documentclass[../../../include/open-logic-section]{subfiles}

\begin{document}

\olfileid[fr]{sth}{ord-arithmetic}{using-addition}
\olsection{Applications de l'addition ordinale}

Grâce à l'addition des ordinaux, nous pouvons calculer
explicitement le rang de divers ensembles, au sens de la
\olref[spine][rank]{defnsetrank} :

\begin{lem}\ollabel{rankcomputation}
Si $\setrank{A} = \alpha$ et $\setrank{B} = \beta$, alors :
\begin{enumerate}
	\item\ollabel{exrankpow}
	$\setrank{\Pow{A}} = \alpha\ordplus 1$
	\item\ollabel{exrankpair}
	$\setrank{\{A, B\}} = \max(\alpha,\beta) \ordplus 1$
	\item\ollabel{exrankcup}
	$\setrank{A \cup B} = \max(\alpha,\beta)$
	\item\ollabel{exranktuple}
	$\setrank{\tuple{A,B}} = \max(\alpha,\beta) \ordplus 2$
	\item\ollabel{exranktimes}
	$\setrank{A \times B} \leq \max(\alpha,\beta)
	\ordplus 2$
	\item\ollabel{exrankunion}
	$\setrank{\bigcup A} = \alpha$ si $\alpha$ est
	vide ou limite ; $\setrank{\bigcup A} = \gamma$
	si $\alpha = \gamma\ordplus 1$
\end{enumerate}
\end{lem}

\begin{proof}
Nous invoquerons à plusieurs reprises
la \olref[spine][rank]{ranksupstrict}.

\emph{\olref{exrankpow}.} Si $x \subseteq A$, alors
$\setrank{x} \leq \setrank{A}$. Donc
$\setrank{\Pow{A}} \leq \alpha \ordplus 1$.
Comme, en particulier, $A \in \Pow{A}$,
$\setrank{\Pow{A}} = \alpha \ordplus 1$.

\emph{\olref{exrankpair}.} Par la
\olref[spine][rank]{ranksupstrict}.

\emph{\olref{exrankcup}.} Par la
\olref[spine][rank]{ranksupstrict}.

\emph{\olref{exranktuple}.} Par
\olref{exrankpair}, deux fois.

\emph{\olref{exranktimes}.} Remarquons que
$A \times B \subseteq \Pow{\Pow{A \cup B}}$,
et invoquons \olref{exranktuple}.

\emph{\olref{exrankunion}.} Si
$\alpha = \gamma\ordplus 1$, il existe
$c \in A$ tel que $\setrank{c} = \gamma$,
et aucun !!{element} de $A$ n'a un rang supérieur ;
ainsi, $\setrank{\bigcup A} = \gamma$.
Si $\alpha$ est un ordinal limite, les !!{element}s
de $A$ ont des rangs arbitrairement proches de
$\alpha$, mais strictement inférieurs à lui ;
il en va donc de même des !!{element}s de
$\bigcup A$, d'où
$\setrank{\bigcup A} = \alpha$.
\end{proof}
\noindent
Nous laissons en exercice l'explication de
l'\emph{inégalité} dans \olref{exranktimes}.

\begin{prob}
Trouvez des ensembles $A$ et $B$ tels que
$\setrank{A \times B} =
\max(\setrank{A}, \setrank{B})$.
Trouvez des ensembles $A$ et $B$ tels que
$\setrank{A \times B} =
\max(\setrank{A}, \setrank{B}) \ordplus 2$.
\footnote{Note de l'édition française : le signe
d'égalité manque entre les deux expressions dans
la deuxième demande de l'original ; il est rétabli
ici.} D'autres valeurs du rang sont-elles possibles ?
\end{prob}

Nous pouvons aussi montrer à présent que plusieurs
façons raisonnables de dire qu'un ordinal est «~fini~»
ou «~infini~» coïncident :

\begin{lem}\ollabel{ordinfinitycharacter}
Pour tout ordinal $\alpha$, les propriétés suivantes
sont équivalentes :
\begin{enumerate}
	\item\ollabel{ord:notinomega}
	$\alpha\notin \omega$, c'est-à-dire que
	$\alpha$ n'est pas un nombre naturel
	\item\ollabel{ord:omegaplus} $\omega \leq \alpha$
	\item\ollabel{ord:oneplus}
	$1 \ordplus \alpha = \alpha$
	% \alpha \approx \alpha \ordplus 1$
	\item\ollabel{ord:plusone}
	$\alpha \approx \alpha\ordplus 1$,
	c'est-à-dire que $\alpha$ et
	$\alpha\ordplus 1$ sont équipotents
	\item\ollabel{ord:infinite}
	$\alpha$ est infini au sens de Dedekind
\end{enumerate}
\end{lem}
\noindent
Nous avons ainsi cinq caractérisations équivalentes
de la (non-)finitude d'un ordinal.

\begin{proof}
\emph{\olref{ord:notinomega} $\Rightarrow$
\olref{ord:omegaplus}.} Par trichotomie.

\emph{\olref{ord:omegaplus} $\Rightarrow$
\olref{ord:oneplus}.} Fixons
$\alpha \geq \omega$. Par récurrence transfinie,
il existe un plus petit ordinal $\gamma$
(éventuellement $0$) tel que
$\alpha = \beta \ordplus \gamma$
pour un ordinal limite $\beta$. Alors :
\[
	1 \ordplus \alpha =
	1 \ordplus (\beta \ordplus \gamma) =
	(1 \ordplus \beta) \ordplus \gamma =
	\supstrict_{\delta < \beta}
	(1 \ordplus \delta) \ordplus \gamma =
	\beta \ordplus \gamma =
	\alpha.
\]
\emph{\olref{ord:oneplus} $\Rightarrow$
\olref{ord:plusone}.} Il existe manifestement
!!a{bijection}
$f \colon (\alpha \disjointsum 1) \to
(1 \disjointsum \alpha)$. Si
$1 \ordplus \alpha = \alpha$, il existe un
isomorphisme
$g \colon (1 \disjointsum \alpha) \to \alpha$.
Considérons alors $\comp{f}{g}$.

\emph{\olref{ord:plusone} $\Rightarrow$
\olref{ord:infinite}.} Si
$\alpha \approx \alpha \ordplus 1$,
il existe !!a{bijection}
$f \colon (\alpha \disjointsum 1) \to \alpha$.
Définissons $g(\gamma) = f(\gamma, 0)$
pour tout $\gamma < \alpha$ ; cette
!!{injection} témoigne que $\alpha$ est
infini au sens de Dedekind, puisque
$f(0,1) \in \alpha \setminus \ran{g}$.

\emph{\olref{ord:infinite} $\Rightarrow$
\olref{ord:notinomega}.} C'est
la \olref[z][infinity-again]{naturalnumbersarentinfinite}.
\end{proof}

\end{document}
